\chapter{Decisões: o que levar das outras VMs}
\label{cap:decisoes-vm}

\proximasaulas

Os capítulos anteriores compararam VMs. Este fecha a Parte~IV com a
pergunta prática: o que muda no código e na arquitetura de um serviço
Java/Quarkus por causa disso? Mesmo formato da Parte~III
(capítulo~\ref{cap:decisoes-quarkus}): quando escolher, o que medir, o
que não vem junto.

\section{Quando a BEAM seria a ferramenta melhor}

Pelas propriedades vistas no capítulo~\ref{cap:comparacao-vms}, a BEAM
leva vantagem de projeto quando o sistema tem estas características:
\begin{itemize}
\item \emph{Muitas atividades longas, independentes e com estado
  próprio}: conexões, sessões, dispositivos. Um processo por atividade é
  barato (327 palavras ao nascer) e o limite padrão passa de um milhão
  por nó (seção~\ref{sec:vms-concorrencia}).
\item \emph{Falha parcial como rotina}: o sistema precisa continuar
  quando uma parte quebra. Foi o problema da tese de Armstrong, que cita
  o AXD301 da Ericsson, com mais de um milhão de linhas de Erlang, como
  \enquote{one of the most reliable products ever made by Ericsson}
  (palavras do autor, no resumo da tese).
\item \emph{Justiça de latência entre tarefas que misturam CPU e espera}:
  a preempção por reductions impede que uma tarefa de CPU segure um
  scheduler (com a exceção das NIFs).
\item \emph{Atualizar sem derrubar conexões}: troca de módulo a quente,
  com duas versões convivendo (seção~\ref{sec:vms-hot-code}).
\end{itemize}
Cada vantagem tem preço: mensagens copiadas, NIFs que bloqueiam
schedulers, um JIT que traduz instrução por instrução, sem especular
(seção~\ref{sec:vms-jit-aot}). Não encontramos benchmark publicado e
rigoroso que compare BEAM e HotSpot num serviço como os da turma; por
isso este capítulo não diz qual é \enquote{mais rápida}.

\section{Por que não estamos nela}

Porque a VM vem junto com linguagem, bibliotecas, ferramentas de
diagnóstico e o que a equipe sabe operar. Os serviços da turma são Java
com Quarkus, e a Parte~II mostrou o que se ganha conhecendo a HotSpot a
fundo: JFR, \texttt{jcmd}, JIT, GC. E a JVM fechou parte da distância: o
JEP~444 apresenta as virtual threads como o mesmo tipo de \emph{user-mode
thread} que os processos do Erlang (seção~\ref{sec:vms-concorrencia}).
O que a JVM não dá é o isolamento garantido pela VM. Esse você constrói
com desenho, e é disso que trata a próxima seção.

\section{O que trazer da BEAM para Java e Quarkus}

\footnotesize
\begin{longtable}{@{}>{\raggedright\arraybackslash}p{3cm} >{\raggedright\arraybackslash}p{4.6cm} >{\raggedright\arraybackslash}p{4.6cm}@{}}
\toprule
Na BEAM & Em Java/Quarkus & O que não vem junto \\
\midrule \endhead
Um processo por atividade & Uma virtual thread por tarefa, sem pool
  (seção~\ref{sec:virtual-threads}) & Heap compartilhado; sem preempção
  por tempo \\
\addlinespace
Nada compartilhado; mensagens copiadas & \texttt{record}s imutáveis,
  nenhum estado mutável \texttt{static}, eventos entre componentes &
  A VM não impede compartilhar: é disciplina de código e de revisão \\
\addlinespace
\emph{Let it crash} & Tratar só a exceção prevista; o resto sobe até a
  borda e a requisição falha inteira & O estado que a thread já alterou
  no heap continua lá \\
\addlinespace
Supervisor e árvore de supervisão & \texttt{StructuredTaskScope}
  (\citetitle{doc:jep505}, prévia no JDK~25); \emph{liveness probe} e
  reinício do container (\citetitle{doc:k8s-probes}) & A API é prévia;
  o reinício pelo Kubernetes é do container inteiro \\
\addlinespace
Limite de reinícios do supervisor & Circuit breaker
  (seção~\ref{sec:resiliencia-quarkus}): para de insistir depois de
  falhas seguidas & Protege a dependência, não reinicia nada \\
\addlinespace
Isolamento de falha entre processos & Bulkhead e \texttt{Semaphore}
  por recurso escasso & Isola recursos (threads, conexões), não memória \\
\addlinespace
Troca de código a quente & Deploy sem downtime com contratos
  compatíveis entre versões (seção~\ref{sec:vms-hot-code}) & Duas
  versões convivem em réplicas diferentes, não na mesma VM \\
\bottomrule
\end{longtable}
\normalsize

\paragraph{\emph{Let it crash} numa requisição Java.}
A distinção de Armstrong entre exceção (prevista, com tratamento
conhecido) e erro (ninguém sabe o que fazer) vira uma regra de código:

\begin{lstlisting}[language=Java]
// Defensivo: engole qualquer erro e segue com estado duvidoso.
try { pedido.aplicarCupom(cupom); }
catch (Exception e) { log.warn("ignorando cupom", e); }
repositorio.salvar(pedido);

// Prevista: cupom inválido tem resposta de negócio (ex.: HTTP 422).
// Não prevista: um bug sobe até a borda e a requisição falha inteira,
// sem salvar um pedido pela metade.
try { pedido.aplicarCupom(cupom); }
catch (CupomInvalido e) { throw new RequisicaoInvalida(e.getMessage()); }
repositorio.salvar(pedido);
\end{lstlisting}

O segundo trecho só é seguro se \texttt{pedido} não for compartilhado com
outra thread e se a transação for desfeita quando a exceção subir. É o
\enquote{processo leva o estado junto} do Erlang, feito à mão: objeto
por requisição, transação por requisição.

\section{Quando escolher nativo}

A seção~\ref{sec:jvm-vs-nativo} deu os números e a regra do guia do
Quarkus (comece em JVM; vá para nativo com uma necessidade concreta). A
Parte~IV acrescenta o porquê de cada lado:
\begin{description}
\item[O nativo ganha] quando startup e memória por réplica dominam o
  custo: o trabalho foi feito no build e o processo começa com o heap
  pronto (seção~\ref{sec:vms-jit-aot}).
\item[O nativo perde] quando o serviço vive muito e o throughput de pico
  importa: sem perfil e sem desotimização, não há a especulação herdada
  do Self (seção~\ref{sec:hist-self}); e quando o código depende de
  carregar classes em runtime, porque o mundo é fechado
  (seção~\ref{sec:vms-hot-code}).
\item[Confira o GC] O Serial GC é o padrão do Native Image, e o G1 só
  existe no Oracle GraalVM em Linux (seção~\ref{sec:vms-gc}). Meça o p99
  sob carga antes de decidir.
\item[Meça] Com a metodologia do capítulo~\ref{cap:benchmarks}: startup
  com uma VM nova por medição, steady state com várias iterações e
  várias invocações, e intervalo de confiança nos dois.
\end{description}

\section{Onde aprofundar}
\textcite{paper:armstrong-thesis-reliable-systems}, capítulo~4, para a
filosofia de erros; \citetitle{doc:erlang-supervisor} para supervisão;
\citetitle{doc:jep505} e \citetitle{doc:jep444} para o que a JVM tem de
parecido; \citetitle{doc:quarkus-fault-tolerance} e
\citetitle{doc:quarkus-health} para as pontes no Quarkus;
\citetitle{doc:quarkus-building-native} para a decisão de nativo.

\section{Desafios práticos}

\desafio{Procure no seu serviço todos os \texttt{catch (Exception e)}.
Para cada um, decida: é uma exceção prevista, com tratamento de negócio,
ou um erro engolido? Reescreva dois deles no estilo do exemplo
deste capítulo.}

\desafio{Escolha um componente do serviço que guarda estado mutável
compartilhado entre requisições. Redesenhe-o para que o estado pertença a
uma só tarefa e as outras se comuniquem com ele por mensagens (uma fila
ou um \texttt{BlockingQueue}). O que ficou mais simples, e o que ficou
mais caro?}

\desafio{Descreva, para o seu serviço, a \enquote{árvore de supervisão}
que já existe sem você ter desenhado: quem reinicia o quê quando uma
parte falha (circuit breaker, liveness probe, Kubernetes)? Onde falta um
limite de reinícios?}
